\scanpage{11}
\mspage{25}

Enfin, il nous sera commode d'avoir un critère permettant de déduire
de la propriété P une propriété de platitude. Nous utiliserons :

\uline{Lemme} \inserted{4.}3. Soient $A$ un anneau, $(P_i)_{i\in I}$ une famille de $A$-modules,
\inserted{$P$ leur produit} $M$ un $A$-module de présentation finie. Alors l'homomorphisme canonique
\[
M\otimes P\longrightarrow\prod M\otimes P_i
\]
est un isomorphisme.

En effet, écrivant $M$ comme comme un conoyau d'un homomorphisme
de modules libres \inserted{de type fini} $L'\longrightarrow L$, et notant que le foncteur produit est exact
donc exact à gauche, on est ramené au cas où $M=A^n$, $n$ entier $>0$, où
c'est trivial.

\uline{Corollaire} \inserted{4.}4. Supposons $A$ tel que tout module de type fini sur $A$ soit
de présentation finie, par exemple $A$ noethérien. Alors tout produit de
$A$-modules plats est un $A$-module plat.

En effet, il suffit de montrer que $P\otimes$ transforme un monomorphisme
$M\longrightarrow N$ en un monomorphisme, quand $M$ et $N$ sont de type fini; le lemme 3 donne
alors immédiatement le résultat voulu.

\uline{Corollaire} \inserted{4.}5. Supposons que $A$ satisfasse \deleted{\unsure{les hypothèses}}
\inserted{\unsure{la condition}} de \inserted{4.}4. \inserted{par exemple $A$ noethérien} Soit
$P$ un $A$-module \deleted{topologique} isomorphe à \unsure{une} limite projective
\deleted{\unsure{``stricte'' ``dénombr''}}
\deleted{de $A$-modules projectifs de type fini (``stricte'' = par un système projectif strict),
``dénombrable'' par un système proje}
\deleted{Alors $P$ est un $A$}
d'un système projectif strict, \deleted{de $A$}
admettant un système cofinal dénombrable, de $A$-modules projectifs. Alors
$P$ est un $A$-module plat.

En effet, on voit facilement que sous les conditions dites, $P$ est
isomorphe à un produit de $A$-modules projectifs, donc plats, d'où le résultat.

\mspage{10}
